\par\noindent\textbf{(1)} 板间近似均匀电场 \(E=V/d\)，电容 \(C=\varepsilon A/d\)。体场能 \(U_E=\frac12\int\varepsilon E^2d^3r=\varepsilon AV^2/(2d)=CV^2/2\)。在固定板上电荷时应改写为 \(U_E=Q^2d/(2\varepsilon A)\)，两种控制条件对应不同导数。
\par\noindent\textbf{(2)} 无磁场时Maxwell应力张量 \(T_{ij}=\varepsilon(E_iE_j-\delta_{ij}E^2/2)\)。在板法向为场方向的理想平行板中，法向压力 \(p=\varepsilon E^2/2\)，总吸引力幅值 \(F=\varepsilon AV^2/(2d^2)=Q^2/(2\varepsilon A)\)。吸引方向可由闭合曲面积分判定。
\par\noindent\textbf{(3)} 穿过板间任一跨越两导线的面积，位移电流 \[I_D=\frac{d}{dt}\int\boldsymbol D\cdot d\boldsymbol A=\varepsilon A\dot E=\frac{\varepsilon A}{d}\dot V=C\dot V\]。这与导线传导电流连续；若忽略位移电流，安培环路对不同跨面会给矛盾结果。
\par\noindent\textbf{(4)} 固定 \(Q\) 时 \(U_E=Q^2d/(2\varepsilon A)\)，对增大 \(d\) 求导，外界需做功 \(dU=F\,dd\)。固定 \(V\) 时场能 \(CV^2/2\) 随 \(d\) 增大反而下降，但电源交换电荷作功；正确总势是 \(U_E-VQ=-CV^2/2\)，对 \(d\) 求导仍得到吸引力。
